% ============================================================================
%  Guión oral + Preguntas y respuestas — Defensa del seminario MCTS / TSLA
%  Asignatura: Métodos Numéricos · Universitat de València
%  Autor: Carlos Ruiz Navarro
% ============================================================================
\documentclass[11pt,a4paper]{article}

\usepackage[utf8]{inputenc}
\usepackage[spanish]{babel}
\usepackage[a4paper, margin=2.3cm]{geometry}
\usepackage{amsmath, amssymb, amsthm}
\usepackage{mathtools}
\usepackage{xcolor}
\usepackage{enumitem}
\usepackage{titlesec}
\usepackage{hyperref}

\definecolor{uvblue}{RGB}{0,94,184}
\hypersetup{
  colorlinks=false,
  pdfborder={0 0 1},
  allbordercolors={1 0 0},
  pdftitle={Preguntas y respuestas — MCTS TSLA},
  pdfauthor={Carlos Ruiz Navarro}
}

\titleformat{\section}
  {\color{uvblue}\Large\bfseries}{\thesection}{0.6em}{}
\titleformat{\subsection}
  {\color{uvblue!85!black}\large\bfseries}{\thesubsection}{0.5em}{}

\newcommand{\preg}[2]{%
  \par\medskip
  \noindent\textbf{P\,#1.} \emph{#2}\par\smallskip}

\newcommand{\resp}[1]{%
  \noindent #1\par}

\title{\color{uvblue}\bfseries
  Guión oral y preguntas para la defensa\\[2pt]
  \large Monte Carlo Tree Search aplicado a TSLA}
\author{Héctor Orosa, Sofía Ortiz, Carlos Ruiz}
\date{Métodos Numéricos · Universitat de València}

\begin{document}
\maketitle

\noindent
Este documento recoge un banco de preguntas anticipadas para la defensa
del seminario, con respuestas sintéticas. Las preguntas se han clasificado
por temas y marcado por probabilidad estimada de aparición
($\bigstar$~baja, $\bigstar\bigstar$~media, $\bigstar\bigstar\bigstar$~alta).
La audiencia conoce: Monte Carlo clásico (estimador por media muestral, hit/miss,
generación de aleatorios), LLN/TCL y análisis de error (Tema 9 y
\texttt{montecarlo.pdf}), métodos iterativos Jacobi/Gauss--Seidel/SOR
(Tema 6) y normas matriciales/radio espectral (Tema 5). \textbf{No} conoce
MCTS, UCB1, MDP/Bellman ni GBM/Itô.

\tableofcontents
\bigskip

% ============================================================================
\section*{Guión oral — diapositivas 5 a 22 (5 minutos)}
% ============================================================================

\noindent
18 diapositivas en 5 minutos $\approx$ \textbf{17 segundos por slide}.
Una o dos frases por diapositiva, sin pausas entre ellas.
Los tiempos indicados son orientativos.

\newcommand{\slide}[4]{%
  \par\medskip\noindent
  \textbf{[#1 $\cdot$ #2]}\hfill\textit{$\sim$#3\,s}\par\smallskip
  \noindent #4\par}

\slide{5}{Motivación}{20}{%
``Cada sesión bursátil un inversor decide comprar, vender o mantener.
Es un problema de decisión secuencial bajo incertidumbre: el futuro es
desconocido pero podemos asignarle probabilidades.
La pregunta que nos hacemos es si existe una política
sistemáticamente mejor que la estrategia pasiva de comprar y mantener.''}

\slide{6}{Conexión con el Monte Carlo clásico}{25}{%
``La respuesta está en Monte Carlo. Si simulamos $n$ trayectorias de precios
tomando la acción $a$, la media de los retornos converge casi seguro a la
esperanza real por la Ley de los Grandes Números ---
en MCTS, $Q(a)/N(a)$ es exactamente ese estimador.
Además, por el TCL el error decae como $N^{-1/2}$: doblar la precisión
requiere cuadruplicar las simulaciones} %, de ahí las 5.000 iteraciones.''}

\slide{7}{Objetivos}{10}{%
``El trabajo tiene tres objetivos: formalizar el problema como MDP,
implementar el algoritmo MCTS, y compararlo contra Buy \& Hold.''}

\slide{9}{Proceso de Decisión de Markov}{30}{%
``Formalizamos el problema como un Proceso de Decisión de Markov:
cuádrupla de estados, acciones, transiciones y recompensas.
La propiedad de Markov dice que el futuro depende solo del estado actual.
La ecuación de Bellman dice que la acción óptima en cada sesión es la que
maximiza la recompensa inmediata más el valor esperado del estado siguiente
--- pero resolverla directamente es inviable: la distribución de transición
es desconocida y el espacio de estados es continuo.
MCTS la esquiva muestreando esa distribución implícitamente mediante
rollouts GBM: $Q(a)/N(a)$ estima $V^*(s_t)$ para cada acción sin necesidad
de conocer $P$ explícitamente.
Esto ocurre una vez por sesión bursátil --- 506 veces en total.
(La propiedad de Markov significa que para predecir el futuro basta saber
dónde estás ahora y qué decides hacer --- por eso el estado incluye MA20
y RSI: son un resumen del pasado que hace el problema Markoviano.)''}


\slide{10}{El MDP del trading sobre TSLA}{20}{%
``En nuestro caso el estado recoge el precio, la cartera y dos indicadores
técnicos. Hay 11 acciones discretas, de comprar el 100\,\% al efectivo hasta
vender el 100\,\%. La recompensa es el log-retorno relativo a Buy \& Hold ---
si el mercado sube un 5\,\% y el agente solo sube un 3\,\%, la recompensa
es negativa.''}

\slide{11}{Los 4 pasos del MCTS}{25}{%
``El algoritmo repite cuatro pasos:
\textbf{Selección} --- UCB1 baja por el árbol hasta una hoja prometedora.
\textbf{Expansión} --- añade un nodo nuevo con contadores $N$ y $Q$.
\textbf{Rollout} --- simula el futuro con GBM durante 60 días y obtiene un
retorno.
\textbf{Retropropagación} --- sube ese retorno actualizando todos los nodos
del camino.
Tras 5.000 iteraciones se elige la acción más visitada.''}

\slide{12}{Selección: UCB1}{20}{%
``UCB1 suma un término de explotación --- la media de retornos observada ---
y un término de exploración que crece cuando una acción ha sido poco visitada.
Esto garantiza que toda acción se visita infinitas veces, condición necesaria
para aplicar la LLN.
Es análogo a Gauss-Seidel: actualiza una sola componente por iteración.''}

\slide{13}{Rollout: GBM}{15}{%
``Para simular el futuro usamos el Movimiento Browniano Geométrico.
Tiene solución analítica exacta vía Itô: solo hay que muestrear una normal
estándar por paso, sin error de discretización.
Los parámetros se recalibran cada día con los últimos 60 días de datos.''}

\slide{14}{Señales sobre el precio}{15}{%
``Aquí vemos el precio de TSLA con las señales de compra y venta del agente.
En el tramo bajista predominan las ventas; en el alcista, las compras.
La política se adapta al régimen del mercado.''}

\slide{15}{Cartera base 100}{15}{%
``Este es el efecto en la cartera.
MCTS sale al efectivo antes de la corrección del 75\,\% y preserva el
capital, mientras Buy \& Hold cae con el mercado.''}

\slide{16}{Ventaja relativa: Alfa}{15}{%
``El alfa mide el exceso de retorno sobre B\&H día a día.
Pico de casi $+80\,\%$ tras el mínimo de la corrección.
La tendencia es creciente: la ventaja se acumula.''}

\slide{17}{Exposición por intensidad}{10}{%
``Las señales diferenciadas por intensidad revelan que predominan las
operaciones extremas --- comprar o vender el 100\,\%.
Las posiciones intermedias son minoritarias.''}

\slide{18}{Exposición dinámica al activo}{15}{%
``Esto se confirma aquí: el ratio de exposición pasa casi todo el tiempo
en 0 o en 1, nunca en valores intermedios.
MCTS aprendió una política todo-o-nada con exposición media del 60\,\%.''}

\slide{19}{Convergencia UCB1}{15}{%
``Esta gráfica cierra el argumento teórico.
$Q(a)/N(a)$ se estabiliza hacia las 500 iteraciones para cada acción,
y la envolvente decrece como $N^{-1/2}$ --- validando experimentalmente
la cota del estimador de Monte Carlo.''}

\slide{20}{Métricas comparativas}{20}{%
``En números: ROI del 240\,\% frente al 100\,\% de B\&H, y drawdown máximo
del $-10\,\%$ frente al $-75\,\%$.
La tensión honesta: el Sharpe de MCTS es inferior, 1.56 frente a 1.93,
porque la política todo-o-nada genera más volatilidad diaria.''}

\slide{21}{Limitaciones del modelo}{15}{%
``Dos limitaciones.
Primera: el GBM tiene colas ligeras y subestima crashes extremos.
Segunda: los resultados son específicos de este periodo y activo --- falta
validación cruzada.''}

\slide{22}{Conclusiones}{15}{%
``En resumen: $Q/N$ es Monte Carlo clásico, converge por la LLN con error
$\mathcal{O}(N^{-1/2})$.
Y lo más interesante: MCTS aprendió por sí solo una política todo-o-nada,
lo que sugiere que las posiciones intermedias no aportan ventaja bajo una
recompensa relativa a B\&H.''}

% ----------------------------------------------------------------------------
\subsection*{Apéndice: cómo explicar la ecuación de Bellman (diap.\ 9)}
% ----------------------------------------------------------------------------

\noindent
Si alguien pregunta qué representa la ecuación o quieres introducirla
brevemente, usa esta explicación en lenguaje llano:

\begin{itemize}\setlength{\itemsep}{2pt}
  \item $V^*(s)$ --- \emph{cuánto vale estar ahora en el estado $s$ si
        juegas perfectamente a partir de aquí.}
  \item $\max_a$ --- \emph{eliges la mejor acción posible.}
  \item $R(s,a)$ --- \emph{lo que ganas ahora mismo con esa acción.}
  \item $\gamma\sum_{s'}P(s'|s,a)\,V^*(s')$ --- \emph{el valor esperado
        del estado al que llegas, descontado porque el futuro vale un
        poco menos que el presente.}
\end{itemize}

\medskip
\noindent\textbf{Frase para decir en la exposición:}

\medskip
\noindent\textit{``La ecuación de Bellman captura una idea muy intuitiva:
el valor de una situación es la mejor recompensa que puedo obtener ahora
más el valor esperado de la situación a la que llego.
Si consigo resolver esa ecuación para todos los estados a la vez, tengo
la política óptima.
El problema es que en trading los estados son continuos e infinitos ---
no puedo resolverla directamente, y ahí entra MCTS: la estima estado
por estado mediante simulación.''}

\medskip
\noindent\textbf{Si preguntan por qué Bellman da la política óptima:}

\medskip
\noindent\textit{``Es el principio de optimalidad: una política óptima
tiene la propiedad de que, sea cual sea el estado inicial y la primera
decisión, las decisiones restantes deben formar a su vez una política
óptima para el subproblema que queda.
Si $\pi^*$ no maximizara la expresión en algún estado, podría mejorarse
localmente --- contradicción con que era óptima.''}

\bigskip
\noindent\textit{Total estimado: $\approx 300$\,s $= 5$\,min.
Si vas justo, recorta la diapositiva 6 a:
``$Q(a)/N(a)$ es el estimador de Monte Carlo; converge por la LLN con
error $\mathcal{O}(N^{-1/2})$, de ahí las 5.000 iteraciones.''}

\newpage

% ============================================================================
\section{Monte Carlo clásico y convergencia}
% ============================================================================

\preg{1}{¿Qué significa exactamente ``c.s.'' (casi seguro) y por qué no
basta con decir ``tiende a la esperanza''? \hfill$\bigstar\bigstar\bigstar$}
\resp{%
``Casi seguro'' significa
$\mathbb{P}\bigl(\lim_{n\to\infty}\widehat I_n = \mu\bigr)=1$:
es convergencia puntual de la sucesión de variables aleatorias, no solo
en probabilidad. Existe un conjunto de trayectorias de medida nula donde
puede fallar; en el resto, converge. Es la conclusión de la LLN fuerte
(cf.\ \texttt{montecarlo.pdf}).}

\preg{2}{¿Por qué en el estimador del Tema 9 aparece el factor $v$ (volumen)
y aquí parece que no? \hfill$\bigstar\bigstar\bigstar$}
\resp{%
En integración Monte Carlo $\int_D f = v\cdot\mathbb{E}[f(X)]$ con
$X\sim\text{Unif}(D)$. En MCTS estimamos directamente una esperanza
$\mathbb{E}[R_T\mid s,a]$, no una integral con dominio acotado, así que
$v=1$ implícitamente.}

\preg{3}{¿Estás usando la LLN fuerte o la débil? ¿Por qué importa la
distinción? \hfill$\bigstar\bigstar$}
\resp{%
La \textbf{fuerte} (convergencia c.s.). Garantiza que cada trayectoria del
algoritmo converge, no solo ``con probabilidad alta para $n$ grande''.
La débil bastaría para un resultado puntual, pero la fuerte es el
estándar en Monte Carlo.}

\preg{4}{La cota $\mathcal{O}(n^{-1/2})$: ¿asintótica vía TCL o válida para
$n$ finito (Chebyshev/Hoeffding)? \hfill$\bigstar\bigstar$}
\resp{%
Vía TCL es \textbf{asintótica} (intervalo $z_{\alpha/2}\sigma/\sqrt{n}$
para $n$ grande). Para $n$ finito existen cotas tipo Hoeffding/Chebyshev
del mismo orden $n^{-1/2}$ pero con peores constantes. La diapositiva 6
usa la versión TCL.}

\preg{5}{¿Las muestras son realmente i.i.d.? Distintos rollouts comparten
árbol y la selección no es uniforme. \hfill$\bigstar\bigstar$}
\resp{%
Sí, condicionalmente: los rollouts desde una acción \textbf{fija} $a$ en la
raíz son i.i.d., porque el GBM se simula independientemente con los mismos
parámetros $(\hat\mu,\hat\sigma)$. La mezcla entre acciones no lo es,
pero $Q(a)/N(a)$ solo promedia dentro de cada acción.}

\preg{6}{Si el rollout devuelve un retorno acotado, ¿no se obtienen cotas
más fuertes que las del TCL? \hfill$\bigstar$}
\resp{%
Sí: el log-retorno relativo está acotado en la práctica, así que aplicaría
Hoeffding (decay subgaussiano). No lo usamos porque la cota TCL ya es
suficientemente ajustada y conecta con el formalismo del Tema 9.}

% ============================================================================
\section{MCTS y UCB1}
% ============================================================================

\preg{7}{¿De dónde sale el $\sqrt{2}$ en UCB1? ¿Por qué no otro valor?
\hfill$\bigstar\bigstar\bigstar$}
\resp{%
Sale de la demostración de Auer, Cesa-Bianchi \& Fischer (2002): es el valor
que minimiza el regret asintótico bajo Hoeffding cuando los retornos están
en $[0,1]$. Es óptimo solo bajo esa hipótesis; en la práctica se ajusta
como hiperparámetro $c$.}

\preg{8}{¿Por qué eligen \emph{max-visits} y no \emph{max-Q/N}? Si la
verdadera utilidad es la media, ¿no debería ser esa?
\hfill$\bigstar\bigstar\bigstar$}
\resp{%
Una acción con $Q/N$ alto y $N$ pequeño tiene error estándar $\sigma/\sqrt N$
enorme: puede ser ruido. Max-visits es robusta porque UCB1 \emph{solo
visita mucho} aquello que se ha vuelto a confirmar varias veces
(diapositiva~14).}

\preg{9}{¿Qué garantiza que el árbol no crezca sin control? ¿Cuál es el
coste por iteración? \hfill$\bigstar\bigstar$}
\resp{%
El árbol crece $O(\text{iter})$ nodos por decisión (un nodo nuevo por
iteración). Coste por iteración: $O(d)$ con $d$ profundidad de descenso
$+\;O(H)$ del rollout. En este trabajo: $5\,000\times{\sim}60\approx
3\cdot 10^5$ pasos GBM por decisión.}

\preg{10}{Si dos acciones tienen valor verdadero muy parecido, ¿UCB1 las
distingue alguna vez o oscila? \hfill$\bigstar\bigstar$}
\resp{%
UCB1 las visita alternativamente con frecuencia similar; el \emph{regret}
crece $O(\log n)$ pero la decisión final puede oscilar. Es un caso conocido
de acciones \emph{near-optimal}.

\smallskip
\noindent\textbf{¿Qué es el regret?} En teoría de \emph{multi-armed bandits},
el \emph{regret} (arrepentimiento) tras $n$ tiradas es la pérdida acumulada
respecto a haber jugado siempre la mejor acción $a^*$:
\[
  R_n \;=\; n\,\mu^* \;-\; \mathbb{E}\!\left[\sum_{t=1}^{n} X_t\right]
       \;=\; \sum_{a\neq a^*}\!\Delta_a\,\mathbb{E}[N_n(a)],
  \qquad \Delta_a = \mu^* - \mu_a.
\]
Mide cuánto se ha ``perdido'' por explorar acciones subóptimas. UCB1
garantiza $R_n = O(\log n)$, que es el orden óptimo (cota inferior de
Lai--Robbins). Si dos acciones tienen $\Delta_a$ pequeño, el regret sigue
siendo bajo en valor absoluto aunque la frecuencia de ``error'' sea alta:
por eso la decisión final puede oscilar sin que el algoritmo esté
funcionando mal.}

\preg{11}{¿Es UCB1 óptimo? ¿Hay alternativas (UCB-V, Thompson sampling) y
por qué no usarlas? \hfill$\bigstar\bigstar$}
\resp{%
No es óptimo en sentido de regret finito; sí asintóticamente. Alternativas
(Thompson sampling, KL-UCB) suelen tener mejores constantes. No se usan
aquí por simplicidad y por mantener transparente la conexión con la LLN
del Tema 9.}

\preg{12}{En la retropropagación se acumula $Q(s,a)\mathrel{+}=X$ y se
reporta la \textbf{media} $Q/N$. ¿Qué pasaría si en lugar de promediar
guardáramos el \textbf{máximo} retorno observado, $Q(s,a)\leftarrow
\max\{Q(s,a),X\}$, y eligiéramos por max en cada nodo (variante
\emph{minimax} / \emph{best-of}) ? \hfill$\bigstar$}
\resp{%
Esa variante sí se usa, pero en un contexto distinto: el MCTS clásico
para \textbf{juegos deterministas con adversario} (Go, ajedrez, estilo
AlphaGo). Allí no hay azar en las transiciones: el rollout devuelve un
valor reproducible y tiene sentido propagar el máximo (o min/max
alternados, según jugador) porque corresponde a la ecuación de Bellman
determinista
$V^*(s)=\max_a [R(s,a)+V^*(s')]$.

\smallskip
En nuestro problema el entorno es \textbf{estocástico}: el mismo
$(s,a)$ produce retornos diferentes en cada rollout porque el GBM
muestrea $Z\sim\mathcal{N}(0,1)$ distinto cada vez. Si guardáramos el
máximo:
\begin{itemize}\setlength{\itemsep}{0pt}
  \item perdemos la interpretación de $Q/N$ como estimador de Monte Carlo
        por media muestral, y con ella la garantía
        $Q/N \xrightarrow{\text{c.s.}} \mathbb{E}[R_T\mid s,a]$ por la LLN
        (todo el argumento del Tema 9 deja de aplicar);
  \item el estadístico $\max_k X_k$ es \emph{sesgado al alza} y no
        converge a la esperanza, sino al supremo esencial de la
        distribución --- un valor que mezcla señal y ruido extremo;
  \item el agente se volvería ``optimista patológico'': preferiría
        acciones cuya cola derecha es gruesa aunque su retorno medio sea
        peor, lo contrario de lo que queremos en gestión de riesgo.
\end{itemize}

\smallskip
La elección entre \emph{average backup} (lo que hacemos) y \emph{max
backup} es, por tanto, una decisión de modelado guiada por la
naturaleza del entorno: estocástico $\Rightarrow$ media; determinista
con adversario $\Rightarrow$ máximo. Existen híbridos
(\emph{Bellman backups}, \emph{MCTS-Solver}) que combinan ambos cuando
hay nodos terminales seguros, pero quedan fuera del alcance del
trabajo.}

% ============================================================================
\section{Conexión con métodos iterativos (Tema 6)}
% ============================================================================

\preg{13}{La analogía con Gauss--Seidel suena bien, pero aquí no hay sistema
$Ax=b$. ¿Qué juega el papel de $A$ y de $b$?
\hfill$\bigstar\bigstar\bigstar$}
\resp{%
Es analogía, no equivalencia formal. El ``punto fijo'' es
$q^*(a)=\mathbb{E}[R_T\mid s,a]$ y la iteración es
\[
  q^{(k+1)}(a) = (1-\alpha_k)\,q^{(k)}(a) + \alpha_k X_{k+1},
  \qquad \alpha_k = \tfrac{1}{N^{(k)}(a)+1}.
\]
con:
\begin{itemize}\setlength{\itemsep}{0pt}
  \item $a\in\mathcal{A}$ --- acción del MDP cuyo valor estamos estimando.
  \item $q^{(k)}(a)=Q^{(k)}(s,a)/N^{(k)}(s,a)$ --- iterado $k$-ésimo, la
        media muestral de los retornos observados para $a$ tras $k$
        rollouts (análogo del $x^{(k)}$ de Gauss--Seidel).
  \item $q^*(a)=\mathbb{E}[R_T\mid s,a]$ --- valor verdadero de la acción,
        el ``punto fijo'' al que converge $q^{(k)}$ por la LLN
        (análogo de la solución $x^*$ del sistema lineal).
  \item $X_{k+1}$ --- retorno terminal del $(k{+}1)$-ésimo rollout,
        v.a.\ con $\mathbb{E}[X_{k+1}]=q^*(a)$ (análogo estocástico del
        término independiente $b$).
  \item $N^{(k)}(a)$ --- número de visitas a la acción $a$ tras $k$
        iteraciones (contador del MCTS).
  \item $\alpha_k=1/(N^{(k)}(a)+1)\in(0,1]$ --- tamaño de paso de la
        recurrencia; juega el papel del coeficiente de la ``matriz de
        iteración'' (escalar aquí).
\end{itemize}
El ``$b$'' es estocástico ($X_{k+1}$); la ``matriz de iteración'' es escalar
$1-\alpha_k$. No hay $Ax=b$ literal (informe \S3).}

\preg{14}{¿Hay un radio espectral $\rho(Q)<1$ o es analogía suelta?
\hfill$\bigstar\bigstar$}
\resp{%
Aquí ``$Q$'' se reduce a $|1-\alpha_k|<1$, trivialmente cierto porque
$\alpha_k=1/(N+1)\in(0,1]$. No es radio espectral en sentido estricto.}

\preg{15}{¿Por qué Gauss--Seidel y no Jacobi? ¿Qué sería el ``Jacobi del
MCTS''? \hfill$\bigstar\bigstar$}
\resp{%
\textbf{Por qué MCTS+UCB1 es de tipo Gauss--Seidel.}
La distinción clave entre ambos métodos del Tema 6 es el momento en que
se usa cada actualización:
\begin{itemize}\setlength{\itemsep}{0pt}
  \item \emph{Jacobi} actualiza todas las componentes $x_j^{(k+1)}$ usando
        \textbf{solo valores del paso anterior} $x^{(k)}$: las
        componentes se calculan en paralelo y los nuevos valores no se
        ven entre sí hasta el paso $k{+}2$.
  \item \emph{Gauss--Seidel} actualiza una componente a la vez y, al
        calcular $x_j^{(k+1)}$, usa \textbf{ya los valores recién
        actualizados} $x_1^{(k+1)},\dots,x_{j-1}^{(k+1)}$ (en lugar de
        sus versiones antiguas).
\end{itemize}
MCTS con UCB1 sigue exactamente este patrón sobre el vector
$q^{(k)}=(q^{(k)}(a))_{a\in\mathcal{A}}$:
\begin{enumerate}\setlength{\itemsep}{0pt}
  \item En cada iteración se actualiza \textbf{una sola componente}
        $q^{(k+1)}(a^*)$ --- la de la acción $a^*$ que ha elegido UCB1.
        El resto de componentes $q^{(k+1)}(a)=q^{(k)}(a)$ permanecen
        intactas.
  \item En la iteración siguiente, UCB1 elige $a^{**}$ \textbf{ya con la
        información actualizada}: la nueva $Q(a^*)/N(a^*)$ y el nuevo
        $N(a^*)$ entran inmediatamente en el cómputo de la fórmula
        $Q/N + c\sqrt{\ln N(s)/N(a)}$, no en una iteración futura.
\end{enumerate}
La diferencia con Gauss--Seidel clásico es que el orden de actualización
no es cíclico fijo ($a_1, a_2, \dots, a_{|\mathcal{A}|}, a_1, \dots$) sino
\textbf{adaptativo}: UCB1 decide en cada paso qué componente conviene
refinar en función de la información acumulada. Por eso decimos que
MCTS+UCB1 es ``Gauss--Seidel adaptativo'' (diapositiva~10).

\smallskip
\textbf{Qué sería el ``Jacobi del MCTS''.}
Lanzar en cada iteración un rollout \textbf{para cada acción} en bloque
y actualizar todas las $q^{(k+1)}(a)$ a la vez con los retornos del paso
$k$. Existe en la literatura (\emph{root parallelization}): es trivial
de paralelizar, pero no aprovecha la información cruzada que aporta UCB1
(no concentra el presupuesto de iteraciones en las acciones más
prometedoras), por lo que converge peor a igual número de rollouts
totales.}

\preg{16}{El parámetro $c=\sqrt{2}$ de UCB, ¿juega el papel del $\omega$
de SOR? ¿Existe $c^*$ análogo a $\omega^*$? \hfill$\bigstar$}
\resp{%
Sí, juega un papel análogo: $c$ pequeño $\to$ explotación (sub-relajación);
$c$ grande $\to$ exploración (sobre-relajación). Existe un $c^*$ empírico
pero no hay fórmula cerrada como
$\omega^*=2/(1+\sqrt{1-\rho^2})$.}

% ============================================================================
\section{MDP y modelado del problema}
% ============================================================================

\preg{17}{¿De verdad el problema cumple Markov? El precio futuro depende
del histórico (volatilidad, momentum). \hfill$\bigstar\bigstar\bigstar$}
\resp{%
No estrictamente. Lo arreglamos enriqueciendo el estado con $\mathrm{MA}_{20}$
y $\mathrm{RSI}_{14}$, que codifican el pasado reciente. Es Markov
\textbf{respecto al estado aumentado}, no respecto al precio solo.}

\preg{18}{¿Por qué la recompensa es relativa a B\&H y no log-retorno
absoluto? \hfill$\bigstar\bigstar\bigstar$}
\resp{%
Para que el agente compita explícitamente con B\&H. Con recompensa absoluta,
``no hacer nada'' en mercado alcista parecería buena (sube). La relativa
penaliza el coste de oportunidad (diapositiva~8 / informe~\S4).}

\preg{19}{¿Por qué solo 11 acciones discretas y no espacio continuo?
\hfill$\bigstar\bigstar$}
\resp{%
MCTS necesita árbol finito para que UCB1 visite cada rama infinitas veces.
Espacio continuo requeriría variantes (HOO, \emph{progressive widening})
fuera del alcance del trabajo.}

\preg{20}{El estado incluye RSI y MA20: ¿no es \emph{feature engineering}
arbitrario? \hfill$\bigstar\bigstar$}
\resp{%
Sí. Sin ellos la política sigue funcionando peor pero supera a B\&H;
el trabajo no incluye \emph{ablation} sistemático (limitación L2).}

\preg{21}{¿No deberíais descontar ($\gamma$) las recompensas futuras?
\hfill$\bigstar$}
\resp{%
Horizonte fijo $H=60$ días con $\gamma=1$ es estándar para problemas
episódicos cortos. La recompensa es solo terminal, así que $\gamma$ solo
cambiaría escala.}

% ============================================================================
\section{GBM y modelo de precios}
% ============================================================================

\preg{22}{El GBM supone log-retornos gaussianos, pero los reales tienen
colas pesadas. ¿Hasta qué punto invalida los rollouts?
\hfill$\bigstar\bigstar\bigstar$}
\resp{%
Es la limitación L1: el GBM infraestima crashes. En el backtest la
corrección $-75\%$ ocurrió pese al modelo, y el agente la gestiona porque
la calibración rolling sube $\hat\sigma$ y tira a efectivo. La garantía
teórica del rollout sí queda comprometida.}

\preg{23}{¿Por qué habéis usado GBM y no otro modelo de precios?
\hfill$\bigstar\bigstar\bigstar$}
\resp{%
El GBM es el modelo \textbf{más simple} que cumple las tres propiedades
mínimas que necesita el rollout del MCTS:

\begin{itemize}\setlength{\itemsep}{0pt}
  \item \textbf{Precios siempre positivos.} Al partir de
        $dS_t = \mu S_t\,dt + \sigma S_t\,dW_t$, la solución
        $S_{t+h}=S_t\exp(\cdots)$ es estrictamente positiva por
        construcción --- coherente con la realidad del mercado.
  \item \textbf{Log-retornos estacionarios y gaussianos.}
        Bajo el modelo, $\ln(S_{t+h}/S_t)\sim
        \mathcal{N}\bigl((\mu-\sigma^2/2)h,\sigma^2 h\bigr)$. Esto
        permite calibrar $\hat\mu$ y $\hat\sigma$ por máxima
        verosimilitud (media y varianza muestrales de los log-retornos)
        en una ventana corta y reciente.
  \item \textbf{Solución analítica exacta vía Itô.}
        $S_{t+h}=S_t\exp((\mu-\sigma^2/2)h+\sigma\sqrt h\,Z)$ con
        $Z\sim\mathcal N(0,1)$. Solo hace falta muestrear un escalar
        $Z$ por paso --- no se integra una SDE (\emph{Stochastic
        Differential Equation}, ecuación diferencial estocástica)
        numéricamente, así que
        \textbf{no hay error de discretización} (a diferencia de
        Euler--Maruyama). Esto es crítico cuando el rollout simula 60
        pasos por iteración y se ejecuta 5\,000 veces por decisión:
        evita que el ruido de discretización contamine la estimación de
        $Q/N$.
\end{itemize}

\smallskip
\textbf{Comparación con alternativas más realistas:}

\begin{itemize}\setlength{\itemsep}{0pt}
  \item \emph{Merton (GBM + saltos de Poisson)}: añade colas pesadas y
        \emph{crashes}. Más fiel a los retornos reales, pero introduce
        $\geq 3$ parámetros adicionales (intensidad y ley de los saltos)
        que requieren ventanas de calibración mucho más largas para ser
        estables.
  \item \emph{Heston (volatilidad estocástica)}: $\sigma_t$ es ella misma
        un proceso. Captura el \emph{volatility clustering} que se
        observa en TSLA, pero es una SDE de dos dimensiones sin solución
        analítica cerrada para el precio: obliga a Euler--Maruyama.
  \item \emph{GARCH(1,1)}: tiempo discreto, modela varianza condicional.
        Es la opción ``natural'' a frecuencia diaria, pero no encaja
        limpiamente con el formalismo del MDP de tiempo continuo que
        usamos.
\end{itemize}

\smallskip
\textbf{Conclusión.} GBM es la elección que minimiza la complejidad del
rollout (un único $Z$ por paso, sin error de discretización) y que mejor
conecta el trabajo con el contenido del curso (lema de Itô y Monte Carlo
del Tema 9). El precio a pagar son las colas ligeras, que reconocemos
como limitación L1 y dejamos como extensión natural en forma de modelo
de Merton.}

\preg{24}{¿Cómo se calibran $\hat\mu$ y $\hat\sigma$? ¿Por qué ventana
de 60 días? \hfill$\bigstar\bigstar$}
\resp{%
Compromiso clásico: corta $\to$ sigue régimen pero ruidosa; larga $\to$
estable pero lenta. 60 días $\approx$ 3 meses, robusto frente a
estacionalidad semanal/mensual. No optimizado.}

\preg{25}{¿$\hat\mu$ con 60 datos no es enormemente ruidoso? El error
estándar de la deriva es brutal. \hfill$\bigstar\bigstar$}
\resp{%
Cierto: error estándar $\hat\sigma/\sqrt{60}\approx \sigma/8$, comparable
a $\mu$. Por eso \emph{Common Random Numbers} con $Z$ compartido es tan
importante: el ruido de $\hat\mu$ se cancela al comparar agente vs.\ B\&H
bajo la misma trayectoria.}

\preg{26}{¿Habéis comprobado normalidad de los log-retornos?
\hfill$\bigstar$}
\resp{%
El informe muestra QQ-plot: colas más pesadas que la normal
(kurtosis~$\approx 5$). Limitación L1.}

\preg{27}{Si la solución es exacta, ¿por qué se llama ``Monte Carlo''?
\hfill$\bigstar$}
\resp{%
Porque lo aleatorio es $Z\sim\mathcal{N}(0,1)$, que es exactamente lo que
se simula. La ``solución exacta'' se refiere a la SDE (no requiere
Euler--Maruyama), no al precio futuro.}

% ============================================================================
\section{Resultados y validez experimental}
% ============================================================================

\preg{28}{$+240\%$ vs.\ $+100\%$: ¿no es sospechoso? ¿No habrá sobreajuste
de hiperparámetros? \hfill$\bigstar\bigstar\bigstar$}
\resp{%
Riesgo real (limitación L2). No hubo \emph{grid search} sobre el período
test, pero los hiperparámetros (5\,000, 60, 60) se fijaron a priori,
no tras ver resultados. Falta evaluación out-of-sample formal.}

\preg{29}{Una sola semilla, un activo, un período. ¿Qué confianza
estadística tiene esto? \hfill$\bigstar\bigstar\bigstar$}
\resp{%
El informe v2 incluye distribución de retornos sobre 100 semillas
(figura \texttt{mcts\_violin\_retornos}) que sí cuantifica varianza,
pero no entró en la presentación. Comparación rigurosa requeriría más
activos y validación cruzada temporal.}

\preg{30}{Sharpe $1{,}56 < 1{,}93$: ¿no es ese el ratio que importa al
inversor? \hfill$\bigstar\bigstar$}
\resp{%
Honestamente, sí. Defensa: el drawdown $-10\%$ es preferible al $-75\%$
para un inversor real con horizonte vital finito (Sharpe penaliza por
igual volatilidad al alza y a la baja). \textbf{Sortino} sería más justo.}

\preg{31}{¿Otros benchmarks (rebalanceo, momentum simple)?
\hfill$\bigstar\bigstar$}
\resp{%
No probados. Limitación reconocida; B\&H es el benchmark estándar pero
laxo.}

\preg{32}{¿Costes de transacción, slippage, impuestos? \hfill$\bigstar\bigstar$}
\resp{%
No incluidos. Una política bang-bang con muchas operaciones grandes
implica que con comisiones realistas ($\sim$5~bps) el alfa caería
notablemente. Discutido en informe~\S6.}

\preg{33}{¿La política reacciona o ``huele'' la corrección por la
calibración? \hfill$\bigstar$}
\resp{%
El GBM rolling captura cambios de régimen (vol sube $\to$ rollouts más
dispersos $\to$ MCTS prefiere efectivo). No es información futura, pero
le da reflejos. Reproducible con otras semillas.}

% ============================================================================
\section{Implementación y reproducibilidad}
% ============================================================================

\preg{34}{¿Cuánto tarda una decisión? \hfill$\bigstar\bigstar$}
\resp{%
${\sim}2$--$3$~s en CPU estándar para 5\,000 iter $\times$ 60 días.
Backtest completo (506 días): ${\sim}25$~min.}

\preg{35}{¿Cómo se paraleliza? \hfill$\bigstar\bigstar$}
\resp{%
Trivial a nivel de rollouts (\emph{root-parallel}, \emph{leaf-parallel}).
No implementado: ejecución secuencial vectorizada con NumPy.}

\preg{36}{¿Por qué Python y no Cython/C++? \hfill$\bigstar$}
\resp{%
Suficiente: los 60 pasos GBM están vectorizados con NumPy. Cuello de
botella: descenso por el árbol (Python puro). Cython aceleraría
${\sim}5\times$.}

\preg{37}{¿La semilla 42 cambia mucho los resultados frente a otras?
\hfill$\bigstar$}
\resp{%
Sobre 100 semillas (informe v2): ROI medio $+180\%$, desviación
$\pm 60\%$. La semilla 42 está cerca de la mediana; no es \emph{cherry
picking} extremo.}

% ============================================================================
\section{Teoría más profunda}
% ============================================================================

\preg{38}{¿La política aprendida es $\pi^*$ óptima o solo casi-óptima?
¿Hay teorema de optimalidad? \hfill$\bigstar\bigstar$}
\resp{%
Casi-óptima asintóticamente: Kocsis \& Szepesvári (2006) prueban que el
regret de UCT es $O(\log n)$ y la política max-visits converge a la
óptima del MDP. Para $n$ finito, no.}

\preg{39}{La convergencia de $Q/N$ es para acción fija desde la raíz.
¿Y la del árbol entero? \hfill$\bigstar\bigstar$}
\resp{%
Mismo resultado de Kocsis--Szepesvári: en cada nodo del árbol,
$Q/N\to V^*$ del subárbol bajo UCT. Citado en el informe~\S3.}

\preg{40}{¿Qué pasa si el MDP no es estacionario (mercado cambia de
régimen)? \hfill$\bigstar$}
\resp{%
Formalmente, en cada decisión el MDP interno es estacionario (porque
$\hat\mu,\hat\sigma$ están fijos para esa decisión). Entre decisiones
cambia, pero cada problema interno se resuelve en su régimen. No hay
garantía global.}

\end{document}
% ============================================================================
